% !TeX root = ../report.tex
\section{Introduction}
\label{sec:introduction}

Mars has a dusty season every year near perihelion, and every observed planet-encircling (global) dust storm has begun in that season. Yet global storms occur in only about a quarter of Mars years, at irregular intervals~\cite{smith2019}. One proposed explanation is that a large storm depletes the loose dust where lifting is easiest, and that dust must be replenished before another can grow~\cite{smith2019,mulholland2013}. This report asks whether a small model can turn a regular dusty season into intermittent large events. The model combines dust-lifting feedback, a reservoir of liftable surface dust that events deplete and that slowly replenishes, and seasonal forcing that varies modestly between years. Simulations and a fixed-dust equilibrium analysis test this question.
